% =====================================================================
% Appendix - prompts e output schemas dos tres bracos de VLM
% Fonte verbatim: paper/pillar23/pillar3-prompts.txt (gerado por dump_prompts.py)
% Validado contra os runs reais e contra forge.vlm.phases / forge.vlm.prepare
% Hashes recalculados em 2026-09-07 importando as funcoes de producao.
%
% Requer no preambulo:
%   \usepackage{listings}
%   \lstset{basicstyle=\ttfamily\scriptsize,breaklines=true,
%           columns=flexible,frame=single,xleftmargin=0pt}
% =====================================================================

\section{VLM prompts and output schemas}
\label{app:prompts}

All three arms judge the same 501 detector proposals over the same 313 reference
frames with the same model, temperature and images. The text below is verbatim;
the placeholders in braces are the only per-proposal substitution. Rendered on
the first proposal of the set (\texttt{20240103201734:4:0}, proposed class
\texttt{person}, detector confidence 0.7559), the prompts hash to:

\begin{center}
\scriptsize
\begin{tabular}{ll}
\toprule
Prompt & SHA-256 \\
\midrule
Single verdict (v1 \emph{and} v3) & \texttt{8295e3a865682eae6c17cdeea17ffb385e2b847dcb2b69fee5a557e0192edbb0} \\
Decomposed (v2)                   & \texttt{ad13f1913a8e46fdf9df1ce84dd39320dbb07f1510cd595dc7b8ee813fbc7d28} \\
\midrule
Schema & SHA-256 (canonical JSON, sorted keys) \\
\midrule
v1 & \texttt{ba932b9570dc4359c681fb10ed81646e2e06404ba1af1e9db1ebf270e4d7fa58} \\
v3 & \texttt{78360f95bcf528e3ed1bc141165a2dee26b88609621383f225688ebb78035197} \\
v2 & \texttt{0eb1eb9349e9af761165e9aef4cb8da76629098e6ff6ca36aecd1cb716557d84} \\
\bottomrule
\end{tabular}
\end{center}

\noindent
Arms v1 and v3 do not merely share prompt text: they call the same function, so
the two prompts are byte-for-byte identical by construction. The SHA-256 over
the concatenation of all 501 rendered single-verdict prompts is
\texttt{b947fef992c369ca3941bc6d46cd3b960e2c1b79ac125de71ceb42feabc57672}.
\textbf{The v1$\to$v3 contrast is therefore a pure schema intervention.}

% ---------------------------------------------------------------------
\subsection{Arm v1 --- single verdict}

\begin{lstlisting}
You are auditing an automatic object detector's output, exactly as a human annotator would. A detector proposed ONE box. Judge only that box.

You get two images: the full driving scene with the proposed box outlined in pink, and a zoomed crop around the same box.

Choose exactly one verdict:

- "kept": the class is right AND the box is tight around the object. No change needed.
- "corrected-minor": the class is right and the box is on the right object, but the edges need a small nudge (slightly loose, slightly cropped).
- "corrected-redrawn": the class is right and the object is there, but the box is badly placed or badly sized - it would have to be drawn again from scratch.
- "class-changed": there IS an object there and the box is reasonable, but the label is wrong (e.g. a truck labelled as a car). Give the correct class.
- "rejected": there is no such object there at all - the detector fired on background, a reflection, a shadow, or a duplicate of another box.

Be decisive. "kept" is a real answer and is common: if the box looks right, say kept rather than hedging toward a correction.

Proposed class: {proposed_class}
Detector confidence: {conf:.2f}
Proposed box (normalized x1,y1,x2,y2): [{x1:.3f}, {y1:.3f}, {x2:.3f}, {y2:.3f}]

Give your verdict and a one-sentence reason.
\end{lstlisting}

\noindent\textbf{Note.} The em dashes in items three and five are U+2014 in the
source; they are rendered as hyphens above for typesetting only.

\begin{lstlisting}[language=,caption={Output schema, arm v1}]
{
  "type": "object",
  "properties": {
    "verdict": {
      "type": "string",
      "enum": ["kept", "corrected-minor", "corrected-redrawn",
               "class-changed", "rejected"]
    },
    "correct_class": {
      "type": "string",
      "description": "Only when verdict == class-changed: the right class."
    },
    "reason": { "type": "string" }
  },
  "required": ["verdict", "reason"],
  "additionalProperties": false
}
\end{lstlisting}

\noindent
In v1, \texttt{correct\_class} is a free-form string and is \emph{absent from}
\texttt{required}. Under constrained decoding the field is therefore optional
and unconstrained, and the model omitted it on all seven occasions it answered
\texttt{class-changed}: the scorer fell back to the original label and credited
zero corrections.

% ---------------------------------------------------------------------
\subsection{Arm v3 --- single verdict with class correction}

\textbf{The prompt text is byte-for-byte identical to arm v1.} The schema is the
only difference:

\begin{lstlisting}[language=,caption={Output schema, arm v3}]
{
  "type": "object",
  "properties": {
    "verdict": {
      "type": "string",
      "enum": ["kept", "corrected-minor", "corrected-redrawn",
               "class-changed", "rejected"]
    },
    "correct_class": {
      "type": "string",
      "enum": ["person", "bicycle", "car", "motorcycle", "bus", "truck"],
      "description": "Only when verdict == class-changed: the right class."
    },
    "reason": { "type": "string" }
  },
  "required": ["verdict", "correct_class", "reason"],
  "additionalProperties": false
}
\end{lstlisting}

\noindent
Exhaustive diff v1 $\to$ v3:

\begin{center}
\small
\begin{tabular}{lll}
\toprule
Field & v1 & v3 \\
\midrule
\texttt{properties.correct\_class.enum} & absent & the six classes \\
\texttt{required} & \texttt{[verdict, reason]} & \texttt{[verdict, correct\_class, reason]} \\
\bottomrule
\end{tabular}
\end{center}

\noindent
Nothing else changes: same five-verdict enum, same property order, same
\texttt{additionalProperties}, same images, same model, same temperature.

% ---------------------------------------------------------------------
\subsection{Arm v2 --- decomposed}

\begin{lstlisting}
You are auditing an automatic object detector, exactly as a human annotator would. A detector proposed ONE box, outlined in pink. You get the full driving scene and a zoomed crop of the same box.

Answer three things, in this order:

1. object_present - Is there really an object of one of these classes inside the outlined box?
   person, bicycle, car, motorcycle, bus, truck
   Answer false ONLY if there is no such object there at all: the detector fired on background, road surface, a reflection, glare, or a shadow.
   If a vehicle IS there, answer true - even if the box is a poor fit. A badly drawn box around a real object is still a real object; you will say so in the next field.

2. object_class - Which of those classes it actually is. Answer this whenever object_present is true. If the detector's guess is right, repeat it; if it is wrong, give the right one.

3. box_quality - How well the box fits that object:
   - tight        : the box hugs the object. No change needed.
   - needs_nudge  : right object, but the edges are slightly loose or slightly cropped.
   - badly_placed : right object, but the box is badly sized or badly positioned and would have to be drawn again.

Be decisive. "tight" is a real answer and is common.

Detector guess: {proposed_class} (confidence {conf:.2f})
Proposed box (normalized x1,y1,x2,y2): [{x1:.3f}, {y1:.3f}, {x2:.3f}, {y2:.3f}]
\end{lstlisting}

\begin{lstlisting}[language=,caption={Output schema, arm v2}]
{
  "type": "object",
  "properties": {
    "object_present": {
      "type": "boolean",
      "description": "Is there really an object of the six classes in the box?"
    },
    "object_class": {
      "type": "string",
      "enum": ["person", "bicycle", "car", "motorcycle", "bus", "truck"]
    },
    "box_quality": {
      "type": "string",
      "enum": ["tight", "needs_nudge", "badly_placed"]
    },
    "reason": { "type": "string" }
  },
  "required": ["object_present", "object_class", "box_quality", "reason"],
  "additionalProperties": false
}
\end{lstlisting}

\noindent
The v2 schema has \emph{no} \texttt{verdict} field. The five-category verdict is
derived in code, with class taking precedence over geometry, matching the
precedence used by the human-verdict accounting:

\begin{lstlisting}[language=,caption={Verdict derivation for arm v2}]
if not object_present:            verdict = rejected
elif object_class != proposed:    verdict = class-changed
elif box_quality == needs_nudge:  verdict = corrected-minor
elif box_quality == badly_placed: verdict = corrected-redrawn
else:                             verdict = kept
\end{lstlisting}

\noindent
Because the constrained decoder builds its grammar in the declared property
order and every property is required, the ordering of the three questions is
enforced by the decoder rather than requested in prose.

% ---------------------------------------------------------------------
\subsection{Enumerations used}

\begin{center}
\small
\begin{tabular}{lll}
\toprule
Enum & Arms & Values \\
\midrule
\texttt{verdict} & v1, v3 & kept, corrected-minor, corrected-redrawn, class-changed, rejected \\
\texttt{correct\_class} & v3 only & person, bicycle, car, motorcycle, bus, truck \\
\texttt{object\_class} & v2 & person, bicycle, car, motorcycle, bus, truck \\
\texttt{box\_quality} & v2 & tight, needs\_nudge, badly\_placed \\
\bottomrule
\end{tabular}
\end{center}

\noindent
\texttt{box\_quality} exists only in v2. In v1 and v3 the same information is
folded into the single \texttt{verdict} enum, where the intermediate category
\texttt{corrected-minor} was never selected --- zero times in either arm, out of
501 proposals each.
